% TOP_PROBLEM: {"id": 4400008, "title": "Pingree open problems — Boyle problem 2", "category_id": 11, "category": {"id": 11, "name": "dynamical_systems", "display_name": "Dynamical Systems", "description": "Problems about long-term behavior of deterministic systems, Hamiltonian dynamics, and periodic orbits.", "slug": "dynamical-systems", "order_index": 11, "created_at": "2026-07-31T15:26:25.671Z"}, "status": "open", "problem_number": "AMR-043-0008", "set_id": 13}
% FIRST_POSTED: 2026-10-01
% ENTRY_KIND: statement_only
% UnsolvedMath ID 4400008; source catalogue code AMR-043-0008
% !TeX program = lualatex
% Definition and sources only; this file is not an LLM solution attempt.
\documentclass[11pt,a4paper]{article}
\usepackage[margin=25mm]{geometry}
\usepackage{fontspec}
\setmainfont{lmroman10-regular.otf}[BoldFont=lmroman10-bold.otf,
 ItalicFont=lmroman10-italic.otf,BoldItalicFont=lmroman10-bolditalic.otf]
\usepackage{amsmath,amssymb,mathtools}
\usepackage{unicode-math}
\setmathfont{latinmodern-math.otf}
\AtBeginDocument{\let\setminus\smallsetminus}
\usepackage{xurl,hyperref}
\hypersetup{colorlinks=true,urlcolor=blue}
\setcounter{secnumdepth}{0}
\setlength{\parindent}{0pt}
\setlength{\parskip}{5pt}
\setlength{\emergencystretch}{3em}
\begin{document}
\section*{Pingree open problems — Boyle problem 2}\label{4400008}
{\small UnsolvedMath ID 4400008 · AMR-043-0008 · Dynamical Systems · AMR Open Problem Lists}
\subsection{Definitions and mathematical statement}
Let $X\subset A^{\mathbb Z}$ and $Y\subset B^{\mathbb Z}$ be nonempty
two-sided subshifts over finite alphabets, with shifts $\sigma$ and $\tau$.
Both carry their product topologies. Assume that $X$ is defined by finitely
many forbidden words and is topologically mixing: for every pair of nonempty
open sets $U,V\subset X$, all sufficiently large $n$ satisfy
$\sigma^n(U)\cap V\ne\varnothing$.

Suppose a homeomorphism $H:X\to Y$ carries each entire shift orbit onto an
entire shift orbit, meaning
\[
 H(\{\sigma^n x:n\in\mathbb Z\})
       =\{\tau^n H(x):n\in\mathbb Z\}\quad\text{for every }x\in X.
\]
Must $Y$ also be a mixing shift of finite type? Both conclusions are part of
the question: a finite forbidden-word presentation must exist for $Y$, and
its shift must be mixing.

The map $H$ is not assumed to commute with the shifts, preserve the direction
of time, or have a continuous integer-valued orbit cocycle. Thus one cannot
transfer a defining local rule through $H$ as through a sliding block
conjugacy. The assumption that the target is a subshift is substantive: it
requires a finite symbolic alphabet and an expansive shift homeomorphism.
The question concerns orbit equivalence of the whole compact systems,
including their periodic points, rather than equivalence after discarding
exceptional or measure-zero subsets.
\subsection{Short English statement}
Does topological orbit equivalence to a mixing finite-type shift force another subshift to be mixing and of finite type?
\subsection{Sources}
\begin{itemize}
\item[S1] ULAM.AI and the UnsolvedMath Contributors, supplied problems.json, record 4400008 (AMR-043-0008), AMR Open Problem Lists. Catalogue identity and source transcription; CC BY 4.0.
\url{https://huggingface.co/datasets/ulamai/UnsolvedMath}.
\item[S2] Open problems from the 3rd Pingree Workshop on Dynamical Systems. (2010), Mike Boyle, Problem 2. Original source indexed by AMR-043-0008.
\url{https://math.huji.ac.il/~mhochman/open-problems/pingree-open-problems.pdf}.
\end{itemize}
\end{document}
